\begin{proposition}[Sixth perturbation axis: non-medium
maintenance stress]
\label{prop:keio-atpm}
Genome-scale flux balance has no temperature or pH state variable;
the standard constraint-based surrogate for temperature-style
stress is escalation of the non-growth-associated ATP maintenance
demand (the \texttt{ATPM} lower bound; defaults $3.15$ iJO1366 and
$6.86$ iML1515~mmol ATP/gDW/h). A pre-screen compared this encoding
head-to-head with proton-leak forcing (the pH-style surrogate): the
maintenance axis spans $29$--$90\%$ wild-type reduction with
four-to-five usable levels in both models, while the proton-leak
axis compresses to two usable levels in iJO1366 (the proton leaves
partly for free with organic-acid secretion) before hard
infeasibility --- ATPM is selected, and the axis carries a
structural property no supply axis has: \emph{zero} wild-type
carbon-sector degeneracy, phosphoglucose-isomerase at-optimum
variability $0.0$ at every level in both models against $45$--$202$
on every supply level. Levels $\{40,60,80,100\}$ (iJO1366) and
$\{60,80,100\}$ (iML1515) drop the wild-type optima from $0.982$ to
$0.695$, $0.496$, $0.298$, $0.099$ and from $0.822$ to $0.398$,
$0.239$, $0.080$. The labels re-stratify in one direction only:
the nine ATP-synthase genes of the \emph{atp} operon are gained
already at the mildest level ($289\to298$ and $286\to295$;
$\kappa=0.981$ on both), and at the $-90\%$ endpoints the gain set
is $42$/$43$ genes --- the \emph{atp} operon, the
cytochrome-$bo_3$ operon (\emph{cyoA}--\emph{cyoD}), the
thirteen-gene NDH-I operon, lower glycolysis (\emph{pgk},
\emph{eno}), \emph{ackA}, \emph{lpd}, and outer-membrane porins
(\emph{ompG}, \emph{ompL}), oxidative phosphorylation enriched
twenty-six-fold over the genome share ($26$ against $3.2$
expected) --- the mirror image of the anaerobic glycolysis
enrichment of Proposition~\ref{prop:keio-o2-limited}; after the
independent-engine adjudication below there are \emph{no} losses at
any level ($\kappa=0.912$ at the endpoints, gains only). The
association splits by arm exactly as the degeneracy reading
predicts, with one refinement. Plain FBA is \emph{strong at
moderate stress} --- $r=+0.949$, $+0.925$, $+0.821$ on iJO1366, all
above the baseline $+0.603$ --- and collapses only at the $-90\%$
endpoints: $r=+0.444$ (iJO1366), $+0.244$ with held-out AUC
$0.550$ (iML1515). Canonical selection restores the endpoint fully
on iJO1366 ($r=+0.937$, AUC $0.991$, MCC $0.974$; moderate levels
$+0.969$, $+0.949$, $+0.964$) and restores the ranking separation
but not the transitive correlation on iML1515 ($r=+0.703$,
$+0.655$, $+0.476$; AUC $0.992$, $0.992$, $0.978$; labels
$\kappa=1.000$ at every level) --- the near-tie boundary of
Remark~\ref{rem:canonical-protocol}. Since the wild-type optimum is
unique at every level on this axis, the plain endpoint collapse
establishes that wild-level flux-variability width is a sufficient
but not necessary indicator: the operative mechanism is per-knockout
vertex instability at near-zero-biomass optima, which canonical
selection repairs wherever the parsimony stage is discriminating.
\end{proposition}

\begin{remark}[What the sixth axis adds]
\label{rem:keio-multiaxis}
The five medium-side supply axes --- carbon source, electron
acceptor, nitrogen, phosphate, iron --- and the non-medium
maintenance-stress axis now bracket the claim from both sides. The
genuine re-stratifications are both regime switches confined to
the module the perturbation rewires: the anaerobic switch gains
the substrate-level-phosphorylation module; the maintenance-demand
switch gains the energy-transduction module, and does so already
at mild stress --- a $40$~mmol/gDW/h ATP demand makes the ATP
synthase itself load-bearing. The nitrogen-substitution losses
remain the only re-wiring-driven losses. The association is the
invariant object under canonical selection on all five supply
axes in both models ($r\ge+0.843$, held-out AUC $\ge0.979$),
is restored on the maintenance axis in iJO1366, and acquires its
measured boundary on the maintenance axis in iML1515 ---
to which the selection subsection now turns.
\end{remark}
